% Síntesis de resultados previos. Sin hipótesis ni deducciones adicionales.
\begin{frame}[label=nav-frame-123]{Síntesis del recorrido teórico}
\small
\begin{enumerate}
  \item \textbf{Identidad de Lie.}
  Comparar los dos cálculos de $\mathcal L_\xi L$ relaciona los momentos:
  $P^{ab}=-2\mathcal R^{ab}$.
  \medskip
  \item \textbf{Criterio de Lanczos--Lovelock.}
  La condición $\nabla_aP^{abcd}=0$ elimina el término
  $-2\nabla^m\nabla^nP_{amnb}$ de la ecuación métrica.
  \medskip
  \item \textbf{Extensión escalar.}
  Los momentos $\Pi^a$ y $F_\phi$ completan la variación y Bianchi.
  El potencial conserva la estructura
  \[
  \mathcal J^{ab}=2P^{abcd}\nabla_c\xi_d
  -4\xi_d\nabla_cP^{abcd},
  \]
  con el momento de curvatura de la teoría completa.
  \medskip
  \item \textbf{Aplicación a $L_0+Q_\beta$.}
  El acoplamiento aporta la divergencia no trivial:
  \[
  \nabla_aP^{attr}=\frac{\ell^2\beta_0p^2}{4r^3}.
  \]
  El modelo no satisface en general la condición LL.
\end{enumerate}
\end{frame}
